\thispagestyle{plain}
\phantomsection
\addcontentsline{toc}{chapter}{ABSTRACT}
\centerline{\zihao{-3}\bfseries ABSTRACT}

\linespread{1.3}\zihao{-4}\bigskip

This report builds an integrated computational framework from the source package CTaylorSeriesSecond20250324, combining high-order Taylor time integration and Newton-type parameter iteration. For high-accuracy simulation of ordinary differential equation systems, we derive recursive expressions of higher-order derivatives and construct a 16th-order local Taylor integrator, then compare it with classical Runge--Kutta solvers in both accuracy and runtime.

For parameter estimation, the identification problem is formulated as a nonlinear least-squares optimization. Newton and Gauss--Newton schemes are developed with explicit objective, gradient and Hessian (or approximate Hessian) structures. Together with derivative-generation modules in the source code, a reusable interface of function-gradient-Hessian evaluation is designed.

Numerical experiments are conducted under different step sizes, noise levels, initial-guess perturbations, and stopping tolerances. Results show that high-order Taylor integration has strong precision advantage for smooth dynamics, while Gauss--Newton is more robust under moderate measurement noise. Damping and adaptive step control further improve convergence when the initial guess is far from the true parameter vector.

\bigskip
\noindent{\bfseries Key words:}
high-order Taylor series; ODE solver; Gauss--Newton method; parameter identification; numerical stability
